\section{Theorem 1 in the life cycle: how far it goes}\label{sec:thm1}

The programme of Section~\ref{sec:prog} rests on Light's Theorem 1, that saving rises with the
interest rate at every state. In the infinite horizon this is open above logarithmic utility. The
life cycle was the reason for hope: a finite backward induction from an explicit terminal
condition, with the per-step loss compounding only $J$ times instead of to a fixed point. This
section reports what the numbers say, at $\beta=0.96$, a seven-state Tauchen chain with
$\rho=0.9$ and $\sigma=0.4$, a zero borrowing limit, $J=60$, and a flat earnings profile
(\texttt{numerics/lc\_check.m}, \texttt{lc\_robust.m}, \texttt{lc\_induct.m},
\texttt{lc\_verify.m}, \texttt{lc\_bound.m}, \texttt{lc\_crit.m}).

\subsection{The statement is its own terminal condition}

The budget constraint makes two statements identical. Since $a'=Ra+y(z)-c$ at either rate,
\[
  g_k(a,z;r_2)\ \ge\ g_k(a,z;r_1)
  \qquad\Longleftrightarrow\qquad
  c_k(a,z;r_2)-c_k(a,z;r_1)\ \le\ (r_2-r_1)\,a ,
\]
so Theorem 1 \emph{is} the affine bound on the rate response of consumption with slope one and
zero intercept. That bound holds with equality in the last period of life, where $c_0=y(z)+Ra$.
The question is therefore whether the terminal condition propagates backward, and it is a
question about one inequality rather than about a pair of constants to be chosen.

\subsection{It propagates, up to a critical risk aversion}

At $\mu=1$ and $\mu=3$ it propagates exactly. Over all sixty stages and the whole wealth region
reachable from zero wealth at birth, the largest violation of $g_k^{(1)}\le g_k^{(2)}$ is zero,
and the consumption bound holds to five parts in $10^{14}$; refining the asset grid from $2000$
to $16000$ points changes nothing. At $\mu=5$ it fails. The failure is a genuine feature of the
household's problem, not of the solution method: the violation is identical at every grid size,
and endogenous-grid and value-function iteration with a pure grid search agree on it to the
resolution of the grid. Saving at zero wealth in the highest income state, in the first period of
a sixty-period life, falls as the rate rises:

\begingroup\small
\begin{center}
\begin{tabular}{@{}lccccc@{}}
\toprule
$r$ & $2\%$ & $3\%$ & $4\%$ & $5\%$ & $6\%$ \\
\midrule
$g(0,z_{\max})$, endogenous grid & $1.7783$ & $1.7638$ & $1.7452$ & $1.7234$ & $1.6993$ \\
$g(0,z_{\max})$, value iteration & $1.7806$ & $1.7606$ & $1.7406$ & $1.7206$ & $1.7006$ \\
\bottomrule
\end{tabular}
\end{center}
\endgroup

Between the two lies a critical degree of risk aversion, and it is sharp. Bisecting on $\mu$ for
the largest violation over all stages gives:

\begingroup\small
\begin{center}
\begin{tabular}{@{}lccc@{}}
\toprule
Horizon $J$ & $r=2\%$ & $r=4\%$ & $r=\lambda=4.17\%$ \\
\midrule
$20$ & $3.79$ & $3.53$ & $3.51$ \\
$40$ & $3.77$ & $3.39$ & $3.36$ \\
$60$ & $3.77$ & $3.39$ & $3.36$ \\
\bottomrule
\end{tabular}
\end{center}
\endgroup

Theorem 1 holds throughout the life cycle below the critical value and fails above it. The value
is stable in the horizon beyond about forty periods, and falls as the rate rises towards the rate
of time preference. Aiyagari's $\mu=3$ is inside it with little room, at $\mu=3.39$ against
$3$; his $\mu=5$ is well outside. So the life cycle buys one of the two open cases and not the
other, and it buys $\mu=3$ only at rates below the rate of time preference, which is where the
equilibrium lies.

\subsection{Why the natural induction does not close}

The backward induction that the transfer lemma suggests compares the two Euler equations at a
state where the lower-rate household saves. Writing $b=g^{(1)}_{k+1}(a,z)$ for the pivot, the
induction hypothesis and monotonicity of $c_k$ give
$c^{(2)}_k(b')\le c^{(1)}_k(b)+\Delta r\,b$ at the relevant points, and the step closes provided
\[
  R_2\,\E\bigl[(c^{(1)}_k(b)+\Delta r\,b)^{-\mu}\bigr]\ \ge\
  R_1\,\E\bigl[c^{(1)}_k(b)^{-\mu}\bigr]\Bigl(1+\frac{\Delta r\,a}{c^{(1)}_{k+1}(a,z)}\Bigr)^{-\mu}.
\]
To first order this asks $\mu\bigl(b/c^{(1)}_k(b)-a/c^{(1)}_{k+1}(a,z)\bigr)\le1/R$: the ratio of
assets to consumption may rise from one age to the next by at most $1/(\mu R)$. It does not. At
zero wealth in the highest income state the left-hand side is large, because such a household
saves a great deal against a certain fall in income, while the right-hand side is about one. The
margin is negative from the third stage of life onward for every degree of risk aversion we tried,
logarithmic utility included, at $-0.12$ for $\mu=1$, $-2.86$ for $\mu=3$ and $-6.29$ for $\mu=5$
per unit of the rate gap. Keeping the term the step discards, the marginal propensity to consume
times the size of the hypothetical violation, bounds the violation by a multiple of $\Delta r$
rather than eliminating it, which is not enough: the aggregate difference the equilibrium argument
must sign is itself of order $\Delta r$.

So the state of the programme is this. Theorem 1 is true in the life cycle for $\mu=3$, and the
statement to prove has the cleanest possible form, its own terminal condition propagated. What is
missing is an argument at zero wealth, where the pointwise Euler comparison is far from holding
and the truth is delivered by the expectation over income states rather than by any single one.
The same zero-wealth state is where the infinite-horizon slack condition was tightest, which
suggests the two problems share a single missing ingredient rather than two.
